\scanpage{35}
\deleted{\unsure{la même}} Soit $u\in\omega$. La condition que $(S,\omega)$ soit un poly comb. est que $\mathbb Z$ opère (via $u$) transitivement sur $S$, \inserted{et $n\geq3$,} ou encore que $S$ est un torseur sous $\mathbb Z_n$ (où $\mathbb Z_n=\mathbb Z/n\mathbb Z$ si $n$ fini, $\mathbb Z_\infty=\mathbb Z$). \deleted{\unsure{Dans la question|Dans l'opération de}} \inserted{Or l'opération de} $\mathbb Z$ sur $S$ déduite de $u^r\in\psi_r(\omega)$ n'est autre que celle qui se factorise par $\mathbb Z_n$ et se déduit de l'op. précédente via l'hom
\diagram{\begin{tikzcd}[ampersand replacement=\&]
\mathbb Z_n \arrow[r,"r\,\mathrm{id}_{\mathbb Z_n}"] \& \mathbb Z_n
\end{tikzcd}}
Donc elle fait de $S$ un $\mathbb Z_n$-torseur ssi $r\,\mathrm{id}_{\mathbb Z_n}$ est un iso, ce qui signifie que $n$ fini et $(r,n)=1$.

(b) est trivial \deleted{\unsure{Si} $(\unsure{r})$} \deleted{On a $\psi_r(\omega)=\psi_{r'}(\omega)$}

\deleted{\unsure{ssi}} (c) $\psi_r(\omega)=\psi_{r'}(\omega)$ signifie
\[
u^{r'}=u^r\quad\text{ou}\quad u^{r'}=(u^{-1})^r\ (=u^{-r})\quad\text{donc}
\]
\[
r'\equiv r\ (n)\quad\text{ou}\quad r'\equiv-r\ (n).
\]

\uline{Remarque} Une structure de $n$-pol. \uline{orienté} sur $S$ équivaut à la donnée d'une permutation circulaire $u$ sur $S$,
% archive sheet number: 34
